\documentclass[10pt,journal,compsoc]{IEEEtran}
\usepackage[utf8]{inputenc}
\usepackage{amsmath,amsfonts,amssymb}
\usepackage{graphicx}
\usepackage{cite}
\usepackage{algorithmic}
\usepackage{textcomp}

\begin{document}

\title{The Epistemology of Single-View 3D Generation: Balancing Information Deficit, Bayesian Priors, and Operational Fidelity}

\IEEEtitleabstractindextext{%
\begin{abstract}
Single-view 3D generation presents a profound intersection of geometric ambiguity, machine learning, and human perception. Previous literature frequently oscillates between treating Large Reconstruction Models (LRMs) as deterministic reconstructors or condemning them as fundamentally incomplete hallucination engines. In this paper, we propose a rigorous, middle-path synthesis. We formalize the non-uniqueness of single-view inversion not as a catastrophic failure, but as a necessary condition solved via Bayesian priors—a process mathematically analogous to human cognitive vision. To move beyond semantic disputes over "high-fidelity," we operationalize fidelity into distinct, computable structural and perceptual metrics. Furthermore, we formalize the Geometric Fragility Index (GFI) using topological distance metrics to render adversarial masking sensitivity experimentally quantifiable. Finally, we reconcile the theoretical incompleteness of these models with their immense economic utility in production rendering pipelines, acknowledging that photorealism is historically an engineering approximation rather than a strict physical truth.
\end{abstract}
}

\maketitle
\IEEEdisplaynontitleabstractindextext
\IEEEpeerreviewmaketitle


\newtheorem{theorem}{Theorem}
\newtheorem{lemma}{Lemma}
\newtheorem{definition}{Definition}

\section{Introduction}
The projection of a 3D environment onto a 2D plane inherently destroys depth information. Consequently, recovering 3D geometry from a single image is mathematically ill-posed. Recent discourse heavily critiques Large Reconstruction Models (LRMs) for "hallucinating" unobserved geometry, suggesting that single-view 3D generation is epistemically unreliable.

However, treating structural uncertainty purely as a system flaw ignores both biological parallels and practical utility. Human vision routinely and successfully navigates 3D spaces by applying cognitive priors to 2D retinal projections [1]. This paper seeks to elevate the analysis of LRMs from philosophical rhetoric to computable, operationalized science. We introduce rigorous mathematical definitions for geometric fragility, decouple definitions of fidelity, and evaluate the pragmatic validity of generative approximations in industrial rendering.

\section{The Information Deficit and Bayesian Parallels}
The fundamental ambiguity of single-view reconstruction is rooted in perspective projection. Let $\Pi: \mathbb{R}^3 \to \mathbb{R}^2$ be the non-injective projection operator. 

\begin{theorem}[Geometric Non-Uniqueness]
For any generic observation $I \in \mathbb{R}^2$, the preimage $\Pi^{-1}(I)$ contains an infinite set of volumetric configurations. Therefore, the deterministic recovery function $f: I \to G_{true}$ does not exist.
\end{theorem}

To converge on a single geometry $\hat{G}$, LRMs sample from a posterior distribution $P(G | I, \Theta)$, where $\Theta$ represents the latent dataset priors. While earlier critiques framed this reliance on priors as a "hallucination," it is mathematically identical to the Bayesian Brain Hypothesis in cognitive science [2]. Just as humans infer the unseen four legs of a chair based on experiential priors, LRMs infer occluded topologies using ShapeNet or Objaverse statistics. Thus, probabilistic completion is not a failure of reconstruction; it is the optimal Bayesian resolution to an underconstrained system.

\section{Operationalizing Fidelity: Structural vs. Perceptual}
The term "high-fidelity" is often attacked as a semantic overreach. The conflict arises from a failure to define the metric space. We propose decoupling fidelity into two operationalized definitions:

\begin{definition}[Structural Fidelity $\Phi_S$]
The geometric alignment between the generated mesh $\hat{G}$ and the unobserved ground truth $G_{true}$. It is quantifiable via the symmetric Chamfer Distance:
\[ \Phi_S = 1 - \left( \frac{1}{|\hat{G}|} \sum_{x \in \hat{G}} \min_{y \in G_{true}} \|x-y\|_2 + \frac{1}{|G_{true}|} \sum_{y \in G_{true}} \min_{x \in \hat{G}} \|y-x\|_2 \right) \]
\end{definition}

\begin{definition}[Perceptual Fidelity $\Phi_P$]
The visual plausibility of the generated mesh when rendered from novel viewpoints, measurable via metrics such as LPIPS [3] between a rendering function $\mathcal{R}(\hat{G}, v)$ and human expectation.
\end{definition}

LRMs optimize heavily for $\Phi_P$, often at the expense of $\Phi_S$. While academic purists may demand $\Phi_S$, the computer graphics industry operates almost exclusively on $\Phi_P$. If a generated tree lacks accurate posterior branching (low $\Phi_S$) but renders beautifully within a forest scene (high $\Phi_P$), the system has achieved its economic and functional objective.

\section{A Computable Geometric Fragility Index}
To empirically study the robustness of these models against masking artifacts and domain shifts, we must formalize adversarial evaluation. We refine the Geometric Fragility Index (GFI) into a computable metric.

Let $M \in \{0,1\}^{H \times W}$ be the ground-truth segmentation mask, and $\epsilon \in \{-1,0,1\}^{H \times W}$ be a discrete masking noise matrix. Let $f(I \odot M)$ represent the mesh generated by the LRM. 

\begin{definition}[Computable GFI]
The Geometric Fragility Index under the Hausdorff distance $d_H$ is defined as:
\[ \text{GFI}(I, M, \epsilon) = \frac{d_H( f(I \odot M), f(I \odot (M \oplus \epsilon)) )}{\|\epsilon\|_1} \]
\end{definition}
By varying $\|\epsilon\|_1$ (the number of perturbed pixels), researchers can computationally plot the fragility gradient of any LRM, moving adversarial testing from conceptual manifesto to empirical benchmark.

\section{Photorealism as an Engineering Approximation}
A persistent critique of LRMs is their failure to perfectly isolate physically-based rendering (PBR) parameters—specifically, their tendency to bake shadows into albedo textures. While this violates strict radiometric constraints, asserting that this precludes "photorealism" relies on a purist definition of graphics.

Historically, industrial photorealism has always relied on approximations. Techniques such as Screen Space Ambient Occlusion (SSAO), baked global illumination, and normal mapping are mathematical cheats designed to emulate physics without the cost of path-tracing [4]. Therefore, while an LRM's baked albedo may be physically incorrect under dynamic relighting, it remains a highly valid and traditional engineering approximation for bounded-view synthesis.

\section{Conclusion}
Single-view 3D generation is characterized by epistemic incompleteness, but this incompleteness does not render the field invalid. By recognizing that LRMs operate as Bayesian inference engines—much like human cognition—we can align our expectations with mathematical reality. By operationalizing metrics like $\Phi_S$, $\Phi_P$, and the GFI, we transition from philosophical absolutism to empirical evaluation. Ultimately, the tension between structural truth and perceptual approximation is not a flaw of LRMs, but the defining paradigm of modern generative computer graphics.

\section*{References}
\begin{enumerate}
    \item Kersten, D., Mamassian, P., \& Yuille, A. "Object Perception as Bayesian Inference", Annual Review of Psychology, 2004.
    \item Friston, K. "The free-energy principle: a rough guide to the brain?", Trends in Cognitive Sciences, 2009.
    \item Zhang, R., et al. "The Unreasonable Effectiveness of Deep Features as a Perceptual Metric", CVPR 2018.
    \item Akenine-Möller, T., Haines, E., \& Hoffman, N. "Real-Time Rendering", A.K. Peters, 2018.
\end{enumerate}


\end{document}
